\begin{table*}[!t]
\caption{Operational interpretation of the main error modes observed in the diagnostic figures.}
\label{tab:error_taxonomy}
\centering
\small
\begin{tabularx}{\textwidth}{@{}l>{\raggedright\arraybackslash}X>{\raggedright\arraybackslash}X@{}}
\toprule
Error mode & Observable symptom & Practical interpretation \\
\midrule
Cloud-edge ramp & Large residuals during rapid adjacent-sample power changes & Weather variables and cloud percentages do not fully encode local cloud motion or cloud-edge enhancement at 5-min cadence. \\
Low-sun transition & Disproportionate relative error near sunrise and sunset & Small solar-time or elevation errors can shift the clear-sky envelope when available power is low. \\
Overcast attenuation & Higher MAE/RMSE under high cloud cover & Total and layer cloud percentages do not measure optical thickness, cloud base, or local irradiance attenuation directly. \\
Seasonal humidity/cloud regime & Month-hour bands with persistent error structure & The tropical coastal site has seasonal regimes that are only partially captured by monthly anomalies and cloud-layer features. \\
Forecast-input mismatch & Sensor-assisted path greatly outperforms weather-input path & Much of the remaining error comes from resource-estimation uncertainty rather than final power-conversion error alone. \\
Nowcast information advantage & 5-min last-value persistence beats the proposed forecast path & Live previous power is an information-rich nowcast feature and should not be confused with previous-day persistence. \\
\bottomrule
\end{tabularx}
\end{table*}
